%% kapa: exam=1 round=30 year=2019 period=1 form=A subject=civil_law subject_ko=민법 seq=35 number=35
%% source: Q-Net 감정평가사 기출문제 (cst00309) artlSeq=5205205 file=extracted/1008/5205205/2019년 제30회 감정평가사 1차시험 문제지/2019년 제30회 감정평가사 1차시험 문제지/2019년도 제30회 감정평가사 1차 1교시 A형.hwp
%% conversion: kordoc 4.24.51 -> tools/gen1.py (auto); review=proofread-sonnet; flags=none
\begin{question}{35}
관습법상 법정지상권에 관한 설명으로 옳지 않은 것은? (다툼이 있으면 판례에 따름)
\begin{choices}
  \item 미등기건물에 대해서는 건물로서의 요건을 갖추었다 하더라도 관습법상 법정지상권이 인정되지 않는다.
  \item 대지와 건물의 소유자가 건물만을 매도하였으나 매수인이 그 건물의 소유를 위하여 매도인과 대지에 관한 임대차계약을 체결하였다면, 특별한 사정이 없는 한 위 매수인은 대지에 관한 관습법상 법정지상권을 포기한 것으로 볼 수 있다.
  \item 건물의 소유를 위한 관습법상 법정지상권을 취득한 자는 이를 취득할 당시의 토지소유자나 이로부터 토지소유권을 전득한 제3자에게 대하여도 등기 없이 그 지상권을 주장할 수 있다.
  \item 관습법상 법정지상권에 기한 대지점유는 정당한 것이므로 불법점유를 전제로 한 손해배상청구는 성립할 여지가 없다.
  \item 가압류 후 본압류 및 강제경매가 이루어지는 경우, 관습법상 법정지상권의 성립요건인 토지와 건물에 대한 소유자의 동일성 판단은 가압류의 효력 발생 시를 기준으로 한다.
\end{choices}
\end{question}
\answer{1}
